\section{Introduction}
\label{sec:intro}

Hyperspectral imaging (HSI) combined with LiDAR provides complementary spectral and spatial information for land-cover classification. In real-world remote sensing, models must continuously incorporate new classes from diverse geographic domains without retraining from scratch or storing past observations---a setting known as exemplar-free class-incremental learning (CIL). While CIL has been extensively studied for natural images~\cite{rebuffi2017icarl,kirkpatrick2017ewc,li2017lwf}, and a few works address HSI-only incremental learning~\cite{deng2024dcprn,xu2025crossacl}, \emph{no prior work tackles cross-domain CIL for joint HSI+LiDAR classification}.

This gap matters because cross-domain incremental learning introduces a compounding challenge: the model must simultaneously handle (i)~class-incremental forgetting, (ii)~distribution shift across geographic domains, and (iii)~modality-specific drift patterns. Standard approaches---either fully freezing the backbone or fully fine-tuning it---both fail in this setting. Freezing preserves old knowledge but cannot adapt to new domains; fine-tuning adapts but catastrophically forgets (EWC degrades to 54.3\% without SHINE normalization (65.9\% with), worse than even a frozen baseline at 59.2\% (72.7\% with SHINE)).

\paragraph{Motivation: Spectral features are inherently more stable.}
Before designing our method, we analyze the cross-domain feature drift of each modality using standard backbones (LSGA-ViT for HSI, CNN for LiDAR), pre-trained on one domain and evaluated on two unseen domains. As shown in Fig.~\ref{fig:drift}, spectral-dominant features (early ViT layers) exhibit minimal drift, while spatial features (deep ViT layers, LiDAR CNN) show 2--7$\times$ larger drift. This pattern holds across all domain pairs, reflecting the physical reality that spectral absorption signatures are material-intrinsic properties invariant to geographic location, while spatial patterns (texture, morphology, elevation) are domain-dependent.

\paragraph{Our approach.}
Guided by this observation, we propose \methodname{}, a modality-asymmetric adaptation strategy for cross-domain multimodal CIL. Our \backbonename{} backbone explicitly separates features into three branches---spectral, HSI-spatial, and LiDAR-spatial---enabling targeted intervention:
\begin{itemize}
    \item The spectral branch is fine-tuned during a brief warm-up (first domain's tasks) and then \emph{frozen} as a stable semantic anchor.
    \item The spatial branches receive per-task low-rank adapters (LoRA) with \emph{asymmetric rank}: LiDAR LoRA rank $= 2\times$ HSI LoRA rank, matching the measured drift ratio.
    \item All task-specific LoRA modules are accumulated at inference, enabling continuous adaptation without task identifiers.
\end{itemize}

\paragraph{Contributions.}
\begin{enumerate}
    \item \textbf{Marathon CIL benchmark:} We establish the first cross-domain, exemplar-free CIL protocol for multimodal remote sensing, spanning 3 geographic domains, 9 incremental tasks, and 32 classes.
    \item \textbf{Drift-motivated asymmetric PEFT:} We provide quantitative evidence that spectral features are more stable across domains, and design a modality-asymmetric low-rank adaptation strategy that exploits this property.
    \item \textbf{Systematic evaluation:} We compare against 9 CIL baselines on the same backbone, test across 3 domain orderings and 2 random seeds, and perform extensive ablation (12 variants), including negative results on orthogonal constraints, spectral routing, and dual-head designs.
\end{enumerate}

\methodname{} achieves 75.7\% average accuracy on the Marathon CIL benchmark, surpassing all exemplar-free methods by +1.1--9.8 percentage points while adding only 36K parameters.
